\documentclass[10pt, letterpaper]{article}
\usepackage{../../../template/template}

% ------------------------
% Impostazioni documento
% ------------------------

\setDocTitle{Manuale Utente}
\setDocVersion{0.1.0}
\setDocData{9/04/2026}
\setDocIndex{Manuale Utente}
\setDocState{1}
\setDocRecipient{2}

\begin{document}

\makeTitlePage

\newpage 
\addChangelog{0.1.0}{9/04/2026}{V. Baleanu}{-}{\noOne}{\firstDraft}
\makeChangelog

\newpage

\tableofcontents

\newpage

\section{Introduzione}

\subsection{Scopo del documento}

\subsection{Glossario}

\subsection{Riferimenti}
\subsubsection{Riferimenti normativi}
\subsubsection{Riferimenti informativi}

\section{Concetti}

\subsection{Tipi di utenti}
\subsubsection{Operatore Sanitario}
\subsubsection{Amministratore}
\subsubsection{Riepilogo degli utenti}

\subsection{Autenticazione e autorizzazione}
\subsubsection{Token}
\subsubsection{Procedura di autenticazione}

\subsection{Dispositivi}
\subsection{Datapoint}
\subsection{Appartamenti}
\subsection{Reparti}
\subsection{Allarmi}
\subsubsection{Regole di allarme}
\subsubsection{Evento di allarme}

\section{Test}
\subsection{Esecuzione dei test di unità}
\subsection{Esecuzione dei test di integrazione}
\subsection{Esecuzione dei test di accettazione}

\section{Requisiti}
\subsection{Browser supportati}
\subsection{Requisiti hardware}
\subsection{Requisiti software}

\section{Accesso e installazione}
\subsection{Utilizzo via web}

\subsection{Installazione dalla \textit{repository} su \textit{GitHub}}
\subsubsection{Preparazione del Sistema}
\paragraph{Installazione di Docker}
\paragraph{Installazione del Software}
\paragraph{Configurazione del file \texttt{.env}}

\subsubsection{Esecuzione del Sistema}
\paragraph{Avvio}
\paragraph{Spegnimento}
\paragraph{Ripristino}

\section{Guida all'utilizzo}
\subsection{Login}
\subsubsection{Primo accesso come Operatore Sanitario}
\subsection{Collegare l'utente Amministratore al relativo \textit{account MyVimar}}
\subsection{Logout}

\subsection{Visualizzare la \textit{Dashboard}}

\subsection{Gestire le notifiche}
\subsubsection{Visualizzare le notifiche}
\subsubsection{Eliminare le notifiche}

\subsection{Gestire i dispositivi di un appartamento}
\subsubsection{Visualizzare i dispositivi}
\subsubsection{Modificare lo stato dei dispositivi}

\subsection{Gestire gli eventi di allarme}
\subsubsection{Visualizzare gli eventi di allarme}
\subsubsection{Risolvere un evento di allarme}
\subsubsection{Visualizzare lo storico degli eventi di allarme}

\subsection{Gestire le regole di allarme}
\subsubsection{Visualizzare le regole di allarme}
\subsubsection{Creare una nuova regola di allarme}
\subsubsection{Abilitare o disabilitare una regola di allarme}
\subsubsection{Modificare una regola di allarme}
\subsubsection{Eliminare una regola di allarme}

\subsection{Visualizzare le \textit{Analytics} di un appartamento}

\subsection{Gestire i reparti}
\subsubsection{Visualizzare tutti i reparti}
\subsubsection{Creare un reparto}
\subsubsection{Assegnare un appartamento ad un reparto}
\subsubsection{Assegnare un Operatore Sanitario ad un reparto}

\subsection{Gestire gli utenti}
\subsubsection{Visualizzare l'elenco degli Operatori Sanitari}
\subsubsection{Creare un nuovo Operatore Sanitario}
% ------------------------
% PARLARE DELLA GENERAZIONE DELLA PASSWORD PER IL PRIMO ACCESSO
% FARE RIFERIMENTO ALLA SEZIONE PRECEDENTE "PRIMO ACCESSO COME OPERATORE SANITARIO"
% ------------------------
\subsubsection{Eliminare un Operatore Sanitario}

\section{Riferimento API}
% LINK A SWAGGER: view4life.me/api/api#

\section{Supporto Tecnico}
Il nostro team è disponibile per offrire supporto in caso di domande o problematiche legate al progetto. 
In presenza di difficoltà tecniche, dubbi sull'utilizzo delle funzionalità o problemi durante l'installazione, 
è possibile contattarci all'indirizzo: snakebyteteam@gmail.com, descrivendo nella maniera più dettagliata possibile il problema.

\end{document}
